% BORRADOR (comentarios del director, 19-oct-2026): reemplaza MainMatter/Cap5/M06-Chapter5.tex. Borrar esta línea al pasarlo.

\pendiente{Cifras provisionales. Todos los juicios de relevancia de este capítulo los asignó el juez local (Llama 3.1 8B), todavía sin la validación humana. Las tablas se reemplazarán con la corrida final; las conclusiones relativas se revisarán con ella.}

Este capítulo presenta los resultados de las dos evaluaciones del trabajo. Las Secciones~\ref{sec:res-pe1} a~\ref{sec:res-pe4} corresponden a la evaluación sumativa: responden las cuatro preguntas de evaluación (Sección~\ref{sec:met-preguntas}) con la configuración que resultó de la Iteración~2. La configuración de la Iteración~1 se aproxima con el sistema denso que usa e5-base-v2 (Sección~\ref{sec:met-sistemas}). La Sección~\ref{sec:res-fallas} corresponde a la evaluación formativa: reúne las fallas que el protocolo reveló durante las dos iteraciones. Salvo que se indique lo contrario, las medias y las pruebas estadísticas usan como unidad la necesidad de información (Sección~\ref{sec:met-estadistica}).

\section{Corpus, consultas y juicios}
\label{sec:res-corpus}

Todas las corridas parten del corpus congelado de la Sección~\ref{sec:met-corpus}: 1.068 ofertas divididas en 4.013 fragmentos. De las 100 consultas de prueba, 85 producen una lista de ofertas (60 en inglés y 25 en español) y corresponden a 60 necesidades de información; las 15 de agregación se evalúan aparte (Sección~\ref{sec:res-agregacion}).

El agrupamiento de profundidad 10 de los nueve sistemas produjo 1.977 pares de necesidad y oferta. El juez interpretó 1.975; los dos restantes cuentan como no relevantes. Asignó la etiqueta 3 al 30,2\,\% de los pares, la 2 al 26,1\,\%, la 1 al 17,0\,\% y la 0 al 26,6\,\%. Cada necesidad tiene en promedio 18,6 ofertas relevantes (etiqueta $\geq 2$), con un mínimo de 3 y un máximo de 36. Estas cifras son mayores que las de la verificación de cobertura de la Sección~\ref{sec:met-consultas}, que con SQL contó solo las ofertas que cumplen todas las condiciones de la necesidad; el juez considera relevantes también las que satisfacen su parte principal aunque falte un requisito secundario (etiqueta 2). Como suelen ser más que las 10 posiciones evaluadas, el Recall@10 de un sistema rara vez llega a 1.

\section{Validación del juez de relevancia}
\label{sec:res-juez}

\pendiente{Kappa binario y ponderado, con IC del 95\,\%, del juez local y del juez final frente a los juicios humanos (401 pares; 103 de calibración y 298 de validación), y matriz de confusión. Si $\kappa \leq 0{,}60$, resultados con el segundo juez (Sección~\ref{sec:met-juicios}).}

\section{Calidad de recuperación (PE1)}
\label{sec:res-pe1}

\subsection{Frente a las líneas base}

La Tabla~\ref{tab:res-pe1-global} resume las métricas de los nueve sistemas.

\begin{table}[htbp!]
    \centering
    \caption{Calidad de recuperación por sistema (media sobre 60 necesidades; juez local)}
    \label{tab:res-pe1-global}
    \begin{tabular}{|l|c|c|c|c|c|}
        \hline
        \textbf{Sistema} & \textbf{P@5} & \textbf{P@10} & \textbf{nDCG@10} & \textbf{MRR} & \textbf{R@10} \\
        \hline
        ILIKE & 0{,}062 & 0{,}054 & 0{,}057 & 0{,}100 & 0{,}026 \\
        \hline
        BM25 & 0{,}595 & 0{,}558 & 0{,}562 & 0{,}729 & 0{,}330 \\
        \hline
        Denso, e5-base-v2 (Iteración~1) & 0{,}695 & 0{,}662 & 0{,}679 & 0{,}861 & 0{,}377 \\
        \hline
        Denso, pgvector exacto & 0{,}712 & 0{,}658 & 0{,}671 & 0{,}829 & 0{,}369 \\
        \hline
        Denso, pgvector HNSW & 0{,}717 & 0{,}661 & 0{,}675 & 0{,}828 & 0{,}371 \\
        \hline
        Denso, Qdrant HNSW & 0{,}712 & 0{,}664 & 0{,}676 & 0{,}822 & 0{,}375 \\
        \hline
        + filtros & 0{,}777 & 0{,}726 & 0{,}726 & 0{,}890 & 0{,}413 \\
        \hline
        + híbrida & 0{,}763 & 0{,}738 & 0{,}724 & 0{,}855 & 0{,}424 \\
        \hline
        + reordenamiento (completa) & 0{,}838 & 0{,}794 & 0{,}802 & 0{,}962 & 0{,}456 \\
        \hline
    \end{tabular}
\end{table}

La búsqueda por subcadena de la plataforma (ILIKE) no devuelve ningún resultado en 78 de las 85 consultas, porque exige que la consulta completa aparezca en el título o en la empresa. Su nDCG@10 de 0,057 ilustra el punto de partida de la plataforma, pero no es una comparación exigente. BM25 devuelve resultados en 84 consultas, y es la comparación principal.

Frente a BM25, la recuperación densa con el modelo multilingüe mejora el nDCG@10 en $+0{,}114$ [$+0{,}059$; $+0{,}170$], y la configuración completa en $+0{,}240$ [$+0{,}174$; $+0{,}307$]. Las 14 comparaciones de nDCG@10 de los sistemas densos frente a las dos líneas base son significativas ($p_{Holm} = 0{,}0014$, el mínimo que permiten 10.000 permutaciones con 14 comparaciones).

Con el modelo multilingüe, la búsqueda aproximada no pierde calidad: HNSW en pgvector frente a la búsqueda exacta, $+0{,}004$ [$-0{,}000$; $+0{,}009$], y Qdrant frente a pgvector, ambos con HNSW, $+0{,}001$ [$-0{,}005$; $+0{,}007$]; ninguna diferencia es significativa.

\subsection{Aporte de cada componente}

La Tabla~\ref{tab:res-escalera} muestra la escalera de componentes: cada fila agrega un componente a la anterior.

\begin{table}[htbp!]
    \centering
    \caption{Escalera de componentes: nDCG@10 y diferencia con el peldaño anterior, con IC bootstrap del 95\,\% (60 necesidades)}
    \label{tab:res-escalera}
    \begin{tabular}{|l|c|c|c|}
        \hline
        \textbf{Sistema} & \textbf{nDCG@10} & \textbf{Diferencia} & \textbf{$p_{Holm}$} \\
        \hline
        Denso, Qdrant HNSW & 0{,}676 & -- & -- \\
        \hline
        + filtros & 0{,}726 & $+0{,}050$ [$+0{,}018$; $+0{,}087$] & 0{,}031 \\
        \hline
        + híbrida & 0{,}724 & $-0{,}002$ [$-0{,}034$; $+0{,}033$] & 1{,}000 \\
        \hline
        + reordenamiento & 0{,}802 & $+0{,}077$ [$+0{,}042$; $+0{,}114$] & 0{,}002 \\
        \hline
        Iteración completa frente al denso & -- & $+0{,}125$ [$+0{,}081$; $+0{,}173$] & 0{,}001 \\
        \hline
    \end{tabular}
\end{table}

Los filtros estructurados y el reordenamiento aportan mejoras significativas; la búsqueda híbrida no. Las métricas secundarias van en la misma dirección: en P@5, los filtros aportan $+0{,}065$ y el reordenamiento $+0{,}075$ ($p = 0{,}016$ y $0{,}014$ sin ajustar), y la búsqueda híbrida $-0{,}013$.

El nDCG@10 de la recuperación densa con e5-base-v2 fue de 0,73 en la primera corrida (13 de octubre) y es de 0,68 en esta. Ambas cifras corresponden al mismo sistema: el mismo modelo, los mismos vectores y los mismos 50 fragmentos agrupados por oferta. La diferencia viene de la medición. En las 53 consultas cuya lista de 10 ofertas es idéntica en las dos corridas, el nDCG@10 baja de 0,806 a 0,724. Dos cambios lo explican. El juez ahora trata como condiciones duras el salario, la modalidad y el tipo de contrato (falla D1 de la Sección~\ref{sec:res-fallas}): de los 1.321 pares juzgados con ambas versiones del prompt, 99 dejaron de ser relevantes y 42 pasaron a serlo. Además, con nueve sistemas en el agrupamiento cada necesidad tiene más ofertas relevantes juzgadas (18,6 frente a 14,0 en promedio), lo que eleva el ordenamiento ideal con que se normaliza el nDCG. Las demás listas difieren porque entre las corridas se reemplazaron dos necesidades sin cobertura en el corpus y se reconstruyó el índice HNSW.

\subsection{Por categoría de consulta}

La Figura~\ref{fig:res-ndcg-categoria} y la Tabla~\ref{tab:res-pe1-categoria} desagregan el nDCG@10 por categoría. Cada categoría tiene 15 necesidades, de modo que estas cifras son descriptivas y no forman parte de las familias de pruebas.

\begin{figure}[htbp!]
    \centering
    \includegraphics[width=0.85\textwidth]{resultados/ndcg-by-category.pdf}
    \caption{nDCG@10 por categoría de consulta (media sobre 15 necesidades por categoría; juez local)}
    \label{fig:res-ndcg-categoria}
\end{figure}

\begin{table}[htbp!]
    \centering
    \caption{nDCG@10 por categoría de consulta (15 necesidades por categoría; juez local)}
    \label{tab:res-pe1-categoria}
    \small
    \begin{tabular}{|l|c|c|c|c|c|c|}
        \hline
        \textbf{Categoría} & \textbf{BM25} & \textbf{e5-base-v2} & \textbf{Denso} & \textbf{+ filtros} & \textbf{+ híbrida} & \textbf{Completa} \\
        \hline
        Exactas & 0{,}753 & 0{,}834 & 0{,}835 & 0{,}835 & 0{,}842 & 0{,}876 \\
        \hline
        Desajuste de vocabulario & 0{,}471 & 0{,}678 & 0{,}651 & 0{,}651 & 0{,}617 & 0{,}764 \\
        \hline
        Con restricciones & 0{,}447 & 0{,}521 & 0{,}583 & 0{,}777 & 0{,}746 & 0{,}832 \\
        \hline
        Ruidosas & 0{,}578 & 0{,}682 & 0{,}638 & 0{,}643 & 0{,}693 & 0{,}735 \\
        \hline
    \end{tabular}
\end{table}

La ventaja de la recuperación densa sobre BM25 es mayor en las consultas de desajuste de vocabulario ($+0{,}180$) y con restricciones ($+0{,}136$) que en las exactas ($+0{,}082$) y las ruidosas ($+0{,}060$). Es el patrón esperado si la ventaja de la recuperación densa proviene de resolver el desajuste de vocabulario que motiva este trabajo (Sección~\ref{sec:problema}). Cada componente actúa donde se espera. Los filtros cambian casi solo las consultas con restricciones ($+0{,}194$; en las demás categorías, como mucho $+0{,}005$). El reordenamiento mejora todas las categorías, sobre todo las de desajuste de vocabulario ($+0{,}147$). La búsqueda híbrida mejora las ruidosas ($+0{,}050$), en las que BM25 recupera términos exactos que la consulta escribe mal o abrevia, pero empeora las de desajuste de vocabulario ($-0{,}034$) y las de restricciones ($-0{,}031$), en las que BM25 aporta candidatas que comparten palabras con la consulta pero no su función.

\subsection{Consultas con restricciones}

La Tabla~\ref{tab:res-restricciones} muestra el cumplimiento de restricciones (Sección~\ref{sec:met-metricas}), que no depende del juez, sobre las 15 necesidades con restricciones.

\begin{table}[htbp!]
    \centering
    \caption{Cumplimiento de restricciones en las 5 y 10 primeras posiciones (15 necesidades)}
    \label{tab:res-restricciones}
    \begin{tabular}{|l|c|c|}
        \hline
        \textbf{Sistema} & \textbf{C@5} & \textbf{C@10} \\
        \hline
        BM25 & 0{,}433 & 0{,}403 \\
        \hline
        Denso, e5-base-v2 (Iteración~1) & 0{,}360 & 0{,}327 \\
        \hline
        Denso, Qdrant HNSW & 0{,}407 & 0{,}373 \\
        \hline
        + filtros & 0{,}953 & 0{,}947 \\
        \hline
        + híbrida & 0{,}953 & 0{,}953 \\
        \hline
        + reordenamiento (completa) & 0{,}973 & 0{,}960 \\
        \hline
    \end{tabular}
\end{table}

Sin filtros, menos de la mitad de las ofertas devueltas cumple las condiciones de la consulta, con cualquier sistema. Con los filtros, el cumplimiento en las cinco primeras posiciones sube de 0,407 a 0,953 ($+0{,}547$ [$+0{,}353$; $+0{,}727$], $p < 0{,}001$ sin ajustar). No llega a 1 porque el clasificador no extrajo el filtro en 2 de las 24 consultas que lo requerían: las dos consultas por prácticas profesionales, en inglés y en español, en las que no reconoció el tipo de contrato. En las otras 22 la extracción fue exacta.

\subsection{Consultas en español y modelo de \textit{embeddings}}

La Tabla~\ref{tab:res-pe1-idioma} compara las 25 consultas en español con sus 25 contrapartes en inglés, que expresan las mismas necesidades.

\begin{table}[htbp!]
    \centering
    \caption{nDCG@10 de las consultas en español y de sus contrapartes en inglés (25 consultas cada una; juez local)}
    \label{tab:res-pe1-idioma}
    \begin{tabular}{|l|c|c|}
        \hline
        \textbf{Sistema} & \textbf{Español} & \textbf{Inglés (mismas necesidades)} \\
        \hline
        BM25 & 0{,}350 & 0{,}593 \\
        \hline
        Denso, e5-base-v2 & 0{,}615 & 0{,}729 \\
        \hline
        Denso, multilingual-e5-base & 0{,}680 & 0{,}697 \\
        \hline
        + filtros & 0{,}741 & 0{,}751 \\
        \hline
        + híbrida & 0{,}693 & 0{,}758 \\
        \hline
        + reordenamiento (completa) & 0{,}795 & 0{,}851 \\
        \hline
    \end{tabular}
\end{table}

BM25 pierde 0,24 de nDCG@10 al pasar del inglés al español, porque las ofertas están en inglés y la consulta no se traduce. El modelo multilingüe casi elimina esa brecha en la recuperación densa (0,680 frente a 0,697). Frente a e5-base-v2, mejora las consultas en español en $+0{,}065$ [$+0{,}007$; $+0{,}128$] ($p = 0{,}046$ sin ajustar; $p_{Holm} = 0{,}28$), empeora levemente las de inglés ($-0{,}023$ [$-0{,}057$; $+0{,}009$] en las 60 consultas en inglés; $-0{,}032$ en las 25 contrapartes de la tabla) y deja igual el total. La búsqueda híbrida es la que más daña al español ($-0{,}048$), coherente con la debilidad de BM25 en ese idioma. La configuración completa alcanza 0,795 en español, frente a 0,350 de BM25. La columna en inglés permite además una comparación más exigente: como las consultas en español son traducciones de las inglesas, BM25 sobre su contraparte en inglés equivale a BM25 con una traducción perfecta de la consulta. Aun frente a esa referencia (0,593), la configuración completa en español la supera en 0,20.

\subsection{Compuerta de sin resultados}

Las métricas anteriores no aplican la compuerta de sin resultados (Sección~\ref{sec:met-sistema}). Con el umbral de 0,134 calibrado en el conjunto de desarrollo, la compuerta se habría activado en 2 de las 85 consultas: \textit{\enquote{protect company systems from hackers}} y una consulta en español que describe a quien construye las pantallas de una web, ambas de desajuste de vocabulario y ambas con ofertas relevantes en el corpus. Su mejor puntaje de reordenamiento fue 0,082. Son el costo real del umbral: dos falsos negativos en consultas que describen la función en lugar de nombrar el cargo.

\subsection{Consultas de agregación}
\label{sec:res-agregacion}

La Tabla~\ref{tab:res-agregacion} muestra la exactitud de la ruta de agregación antes y después de corregir la falla A7 (Sección~\ref{sec:res-fallas}), detectada en esta misma corrida.

\begin{table}[htbp!]
    \centering
    \caption{Exactitud de la ruta de agregación (15 consultas)}
    \label{tab:res-agregacion}
    \begin{tabular}{|l|c|c|}
        \hline
        \textbf{Medida} & \textbf{Antes} & \textbf{Después} \\
        \hline
        Ruta y plantilla esperadas & 11 (0{,}733) & 12 (0{,}800) \\
        \hline
        Respuesta exacta frente a la referencia & 10 (0{,}667) & 11 (0{,}733) \\
        \hline
        Proporción media de filas cubiertas & 0{,}667 & 0{,}734 \\
        \hline
    \end{tabular}
\end{table}

Quedan cuatro respuestas no exactas, todas fallas abiertas (Sección~\ref{sec:res-fallas}):
\begin{itemize}
    \item dos consultas con la palabra \textit{hiring} se enrutan como mixtas (falla A8);
    \item \textit{\enquote{how many jobs pay more than 200k}} elige la plantilla correcta, que lista las ofertas, pero el catálogo no tiene una que las cuente (falla C2);
    \item \textit{\enquote{how many jobs require python}} recibe la plantilla de conteo por tipo de contrato en lugar de la de demanda de habilidades, un error del clasificador.
\end{itemize}

\section{Fidelidad de la generación (PE2)}
\label{sec:res-pe2}

Las 85 consultas de recuperación se respondieron con la ruta real del pipeline, servida desde Qdrant, con el generador a temperatura 0. Las dos consultas en las que actuó la compuerta recibieron la respuesta de sin resultados y no se puntúan. Las métricas de RAGAS se calcularon con el juez local y una ventana de 16.384 \textit{tokens} (falla D2). La Tabla~\ref{tab:res-pe2} muestra los resultados.

\begin{table}[htbp!]
    \centering
    \caption{Fidelidad y relevancia de la respuesta (RAGAS, juez local), por categoría e idioma}
    \label{tab:res-pe2}
    \begin{tabular}{|l|c|c|c|c|}
        \hline
        \textbf{Grupo} & \textbf{Fidelidad} & \textbf{$n$} & \textbf{Relevancia} & \textbf{$n$} \\
        \hline
        Todas & 0{,}679 & 83 & 0{,}751 & 67 \\
        \hline
        Exactas & 0{,}584 & 22 & 0{,}684 & 14 \\
        \hline
        Desajuste de vocabulario & 0{,}758 & 20 & 0{,}731 & 17 \\
        \hline
        Con restricciones & 0{,}730 & 21 & 0{,}914 & 19 \\
        \hline
        Ruidosas & 0{,}650 & 20 & 0{,}643 & 17 \\
        \hline
        Inglés & 0{,}724 & 59 & 0{,}776 & 47 \\
        \hline
        Español & 0{,}568 & 24 & 0{,}690 & 20 \\
        \hline
    \end{tabular}
\end{table}

La relevancia no pudo calcularse en 16 respuestas (19\,\%), porque la salida del juez no tuvo el formato esperado. Las dos métricas son muy dispersas: la mediana de la fidelidad es 0,77, pero 19 respuestas quedan por debajo de 0,5 y 5 obtienen 0. Como en la primera corrida, una parte de esos valores bajos son errores del juez y no del pipeline (falla D3): asigna fidelidad 0 a respuestas cuyas cifras y títulos aparecen literalmente en el contexto. Con este juez, la Tabla~\ref{tab:res-pe2} no es una medida fiable de la fidelidad del pipeline.

La instrucción de responder en el idioma de la pregunta redujo las respuestas en inglés a consultas en español de 21 de 25 en la primera corrida a 4 de 24 (falla A6).

\pendiente{Fidelidad y relevancia con el juez final, y verificación humana de una muestra de unas 30 respuestas.}

\section{Latencia (PE3)}
\label{sec:res-pe3}

La Tabla~\ref{tab:res-pe3-local} compara la latencia de las etapas propias del artefacto en las dos iteraciones, con cada motor. Cada consulta se repitió 30 veces tras 5 de calentamiento. En la Iteración~2 se excluyen las dos consultas en las que actuó la compuerta, que no construyen contexto.

\begin{table}[htbp!]
    \centering
    \caption{Latencia por etapa, mediana y percentil 95 en milisegundos}
    \label{tab:res-pe3-local}
    \small
    \begin{tabular}{|l|c|c|c|c|}
        \hline
        \multirow{2}{*}{\textbf{Etapa}} & \multicolumn{2}{c|}{\textbf{Iteración~1}} & \multicolumn{2}{c|}{\textbf{Iteración~2}} \\
        \cline{2-5}
         & \textbf{Qdrant} & \textbf{pgvector} & \textbf{Qdrant} & \textbf{pgvector} \\
        \hline
        \textit{Embedding} de la consulta & 53,7 / 66,0 & 53,8 / 65,8 & 13,5 / 15,9 & 13,3 / 15,6 \\
        \hline
        Búsqueda densa & 3,7 / 5,2 & 25,3 / 28,0 & 4,7 / 5,9 & 25,8 / 27,4 \\
        \hline
        BM25 & -- & -- & 0,2 / 0,6 & 0,2 / 0,6 \\
        \hline
        Fusión & -- & -- & 0,0 / 0,1 & 0,0 / 0,1 \\
        \hline
        Lectura de las ofertas & 1,2 / 1,8 & 1,1 / 1,5 & 5,3 / 7,6 & 3,2 / 4,1 \\
        \hline
        Reordenamiento & -- & -- & 811,4 / 917,9 & 792,2 / 916,2 \\
        \hline
        Contexto completo & 59,0 / 71,5 & 80,5 / 93,1 & 835,4 / 949,7 & 833,8 / 949,6 \\
        \hline
        Consultas $\times$ repeticiones & \multicolumn{2}{c|}{$85 \times 30$} & \multicolumn{2}{c|}{$83 \times 30$} \\
        \hline
    \end{tabular}
\end{table}

Tres observaciones resumen la tabla:
\begin{itemize}
    \item El \textit{embedding} de la consulta pasó de 54 a 13,5~ms porque en la Iteración~2 el modelo se ejecuta en la GPU; no es un efecto del cambio de modelo.
    \item La búsqueda densa tarda unos 5~ms en Qdrant y unos 26~ms en pgvector, porque el planificador de PostgreSQL elige el barrido secuencial a esta escala (comportamiento M2 de la Sección~\ref{sec:res-motor}). La diferencia, de unos 21~ms, se mantiene entre iteraciones.
    \item El reordenamiento domina la Iteración~2: 811 de los 835~ms del contexto en Qdrant (97\,\%). Con él, la diferencia entre motores desaparece en la latencia total.
\end{itemize}

El reordenamiento es, por tanto, el componente que más aporta a la calidad y el que más cuesta: $+0{,}077$ de nDCG@10 a cambio de unos 0,8~s por consulta en la GPU del equipo de evaluación (Tabla~\ref{tab:art-config}). En la Iteración~1, la generación con el proveedor externo tardaba en la mediana 0,56~s hasta el primer \textit{token} percibido y 1,23~s en total. El reordenamiento agrega a esa espera unos 0,8~s, un costo del mismo orden que la propia generación.

\pendiente{Latencia de la generación en la Iteración~2 (30 consultas $\times$ 3): tiempo hasta el primer \textit{token}, percibido y total.}

\section{Almacenamiento vectorial (PE4)}
\label{sec:res-pe4}

\subsection{A la escala del corpus real}

A la escala del corpus real, el motor no cambia la calidad: los dos devuelven listas equivalentes (Sección~\ref{sec:res-pe1}), y la búsqueda exacta en ambos coincide en el 99,75\,\% de los 10 primeros resultados de 200 consultas con los vectores multilingües (99,85\,\% con los de e5-base-v2; Sección~\ref{sec:art-vectorial}). La Tabla~\ref{tab:res-ann-real} muestra la búsqueda aislada, medida desde la aplicación.

\begin{table}[htbp!]
    \centering
    \caption{\textit{Recall}@10 y latencia de la búsqueda vectorial sobre el corpus real (200 consultas, vectores multilingües)}
    \label{tab:res-ann-real}
    \begin{tabular}{|l|c|c|c|c|}
        \hline
        \multirow{2}{*}{\textbf{Modo}} & \multicolumn{2}{c|}{\textbf{pgvector}} & \multicolumn{2}{c|}{\textbf{Qdrant}} \\
        \cline{2-5}
         & \textbf{\textit{Recall}} & \textbf{p50 (ms)} & \textbf{\textit{Recall}} & \textbf{p50 (ms)} \\
        \hline
        Exacto & 1{,}000 & 24,4 & 0{,}998 & 3,7 \\
        \hline
        Automático (decide el planificador) & 1{,}000 & 23,8 & -- & -- \\
        \hline
        HNSW, $ef_{search} = 16$ & 0{,}918 & 1,9 & 0{,}850 & 2,7 \\
        \hline
        HNSW, $ef_{search} = 64$ & 0{,}990 & 2,3 & 0{,}977 & 2,9 \\
        \hline
        HNSW, $ef_{search} = 256$ & 0{,}998 & 3,1 & 0{,}998 & 3,3 \\
        \hline
    \end{tabular}
\end{table}

Con el índice forzado, pgvector es más rápido que Qdrant a esta escala y tiene más \textit{recall} con el mismo $ef_{search}$. En modo automático, en cambio, su planificador recorre la tabla completa y tarda unos 24~ms, frente a unos 3~ms de Qdrant. Con los vectores multilingües, HNSW pierde más \textit{recall} a $ef_{search}$ bajo que con e5-base-v2 (con $ef_{search} = 16$, 0,918 y 0,850 frente a 0,976 y 0,925), pero desde $ef_{search} = 64$ ambos motores superan 0,97. En el \textit{benchmark}, con $ef_{search} = 64$ y 8 clientes concurrentes, pgvector atiende entre 6.700 y 8.400 consultas por segundo sobre el corpus real, y Qdrant alrededor de 1.000. Con filtros del 1 y el 10\,\%, el planificador de pgvector resuelve la consulta con el índice del atributo y una búsqueda exacta sobre las filas que lo cumplen (\textit{recall} 1,0 en menos de 0,5~ms). Con un filtro del 50\,\% y $ef_{search} = 16$, sin barrido iterativo, solo el 22\,\% de las consultas obtiene la lista completa (comportamiento M1 de la Sección~\ref{sec:res-motor}).

\subsection{Escalamiento con datos sintéticos}

Los conjuntos sintéticos de 100.000 y 500.000 vectores se generaron con la calibración de la Sección~\ref{sec:met-benchmark}. La búsqueda en rejilla eligió una desviación del ruido de 0,5, con un estadístico de Kolmogórov-Smirnov de 0,23 frente a la distribución de similitudes al vecino más cercano del corpus real. Es el mejor ajuste que permite el generador, no un ajuste perfecto: los vecinos sintéticos se distribuyen de forma algo distinta a los reales. Por eso los resultados se interpretan como una comparación entre motores sobre los mismos datos, y no como una predicción del \textit{recall} absoluto que tendría un corpus real de ese tamaño. La Figura~\ref{fig:res-recall-latencia} muestra tres escenarios y la Tabla~\ref{tab:res-bench-sint} resume los puntos con $ef_{search} = 64$.

\begin{figure}[htbp!]
    \centering
    \includegraphics[width=\textwidth]{resultados/recall-latency.pdf}
    \caption{\textit{Recall}@10 frente a la latencia mediana para $ef_{search} \in \{16, \ldots, 512\}$ (mediana de tres corridas). En el corpus real y sin filtro, pgvector usa el índice forzado.}
    \label{fig:res-recall-latencia}
\end{figure}

\begin{table}[htbp!]
    \centering
    \caption{\textit{Benchmark} sintético con $ef_{search} = 64$: \textit{recall}@10, latencia mediana (ms) y completitud (mediana de tres corridas)}
    \label{tab:res-bench-sint}
    \small
    \begin{tabular}{|l|l|c|c|c|}
        \hline
        \textbf{Datos y filtro} & \textbf{Motor} & \textbf{\textit{Recall}} & \textbf{p50} & \textbf{Completitud} \\
        \hline
        \multirow{2}{*}{100.000, sin filtro} & pgvector & 0{,}948 & 0,88 & 1{,}00 \\
         & Qdrant & 0{,}980 & 2,27 & 1{,}00 \\
        \hline
        \multirow{3}{*}{100.000, filtro del 10\,\%} & pgvector & 0{,}603 & 1,00 & 0{,}09 \\
         & pgvector, barrido iterativo & 0{,}914 & 1,81 & 1{,}00 \\
         & Qdrant & 0{,}975 & 1,83 & 1{,}00 \\
        \hline
        \multirow{2}{*}{500.000, sin filtro} & pgvector & 0{,}816 & 1,01 & 1{,}00 \\
         & Qdrant & 0{,}950 & 2,33 & 1{,}00 \\
        \hline
        \multirow{3}{*}{500.000, filtro del 10\,\%} & pgvector & 0{,}531 & 1,00 & 0{,}10 \\
         & pgvector, barrido iterativo & 0{,}840 & 1,37 & 1{,}00 \\
         & Qdrant & 0{,}974 & 2,31 & 1{,}00 \\
        \hline
        \multirow{4}{*}{500.000, filtro del 1\,\%} & pgvector & 0{,}059 & 0,99 & 0{,}00 \\
         & pgvector, barrido iterativo & 0{,}876 & 9,77 & 1{,}00 \\
         & pgvector, búsqueda exacta forzada & 1{,}000 & 9,03 & 1{,}00 \\
         & Qdrant & 0{,}982 & 1,73 & 1{,}00 \\
        \hline
    \end{tabular}
\end{table}

Sin filtro, con el mismo $ef_{search}$, Qdrant obtiene más \textit{recall} a medida que crece la escala (con 500.000 vectores y $ef_{search} = 64$, 0,950 frente a 0,816). A \textit{recall} igualado no hay un ganador claro: con un \textit{recall} de 0,95, pgvector es más rápido ($ef_{search} = 256$, 1,88~ms, frente a Qdrant con $ef_{search} = 64$, 2,33~ms), y con 0,98 lo es Qdrant ($ef_{search} = 128$, 3,30~ms, frente a pgvector con $ef_{search} = 512$, 4,21~ms).

Con filtros, la diferencia es clara. Sin barrido iterativo, pgvector devuelve listas incompletas: con 500.000 vectores, ninguna consulta obtiene la lista completa con un filtro del 1\,\%, y solo el 10\,\% con un filtro del 10\,\%. El barrido iterativo completa las listas, pero con el filtro del 1\,\% el \textit{recall} se estanca entre 0,87 y 0,90 en todo el rango de $ef_{search}$ y la latencia sube a unos 10~ms. Qdrant mantiene listas completas y un \textit{recall} de al menos 0,97 desde $ef_{search} = 64$ con todos los filtros y escalas, por ejemplo 0,982 en 1,73~ms con el filtro del 1\,\% y 500.000 vectores. Con 100.000 vectores y el filtro del 1\,\%, el planificador de pgvector vuelve a elegir la búsqueda exacta sobre las filas filtradas, con \textit{recall} 1,0 en aproximadamente 1~ms.

El estancamiento de pgvector con el filtro del 1\,\% y 500.000 vectores no se debe a los límites del barrido iterativo. Una medición adicional, en las mismas condiciones, aumentó el número máximo de tuplas recorridas (\texttt{hnsw.max\_scan\_tuples}) de 20.000, el valor por defecto, a 100.000 y a 500.000, y la memoria del barrido hasta ocho veces: el \textit{recall} siguió en 0,876. El modo de orden estricto lo empeoró (0,674 con $ef_{search} = 64$). Una explicación compatible es que el barrido se detiene en cuanto reúne las $k$ filas que cumplen el filtro, que son las primeras que encuentra al recorrer el grafo y no necesariamente las más cercanas. En cambio, al impedir el uso del índice HNSW, el planificador resolvió la consulta con el índice del atributo y una búsqueda exacta sobre las 5.000 filas filtradas: \textit{recall} 1,0 en 9,03~ms, la misma latencia que el barrido iterativo. A esta escala y selectividad, el planificador de pgvector elige el índice HNSW cuando la búsqueda exacta sería igual de rápida y exacta; aun así, Qdrant es unas cinco veces más rápido.

En recursos, con 500.000 vectores, pgvector cargó los datos en 16~s y construyó el índice en 133~s; Qdrant tardó 77 y 360~s. pgvector ocupó 4,1~GB de tabla y 2,0~GB de índice, y Qdrant 1,86~GB en total. La memoria pico fue de 7,8~GB en PostgreSQL, con 4~GB de \textit{shared buffers}, y de 4,1~GB en Qdrant. Por esa memoria no se evaluó un conjunto de un millón de vectores (Sección~\ref{sec:met-benchmark}).

Con 8 clientes concurrentes, pgvector atiende más consultas por segundo que Qdrant en casi todos los puntos, por ejemplo unas 2.400 frente a unas 900 con 500.000 vectores sin filtro. Ambas cifras se miden desde clientes en Python, con el protocolo de PostgreSQL y con gRPC respectivamente, sobre una sola máquina, de modo que reflejan también el costo de cada cliente (Sección~\ref{sec:met-validez}).

El \textit{benchmark} se repitió con los vectores de e5-base-v2. Los \textit{recall} fueron entre 0,01 y 0,09 mayores en los mismos puntos, y todas las conclusiones anteriores se mantuvieron.

\section{Fallas reveladas por la evaluación}
\label{sec:res-fallas}

En un sistema RAG muchos defectos son \emph{silenciosos}: el sistema sigue respondiendo con fluidez y de forma plausible, porque el modelo de lenguaje redacta una respuesta con lo que recibe, aunque lo recibido esté incompleto, mal filtrado o mal enrutado. Usando la aplicación, esos defectos no se notan. \citet{Barnett2024} concluyen que la validación de un sistema RAG solo es posible durante su operación. En JobAI, la preparación y la ejecución del protocolo de evaluación revelaron 18 fallas antes de esa etapa, 12 del artefacto y 6 del propio arnés de evaluación, además de dos comportamientos silenciosos del motor de almacenamiento que el diseño debe tener en cuenta. En términos de DSR son resultados de la evaluación \emph{formativa} \citep{Venable2016}: se obtuvieron durante el ciclo de diseño y orientaron las correcciones de las dos iteraciones (Capítulo~\ref{cap:artefacto}).

Las fallas se clasifican en cuatro tipos según su efecto:
\begin{description}
    \item[A. Silenciosas que cambian los resultados.] El sistema devuelve ofertas, cifras o respuestas incorrectas sin ningún error visible.
    \item[B. Silenciosas que cambian el costo.] Los resultados son correctos, pero se obtienen con un costo mayor del necesario.
    \item[C. Visibles.] Rompen una funcionalidad o producen una respuesta evidentemente incorrecta.
    \item[D. Del arnés de evaluación.] No afectan al artefacto, sino a la medición, y habrían invalidado parte de la evaluación.
\end{description}
Los dos comportamientos del motor (M1 y M2) no se cuentan como fallas, porque no son defectos del artefacto; se describen aparte en la Sección~\ref{sec:res-motor}. La Tabla~\ref{tab:res-fallas} las resume. Después se describe cada una con la misma secuencia: observación, causa, efecto, corrección y lección para otros sistemas RAG.

{\small
\begin{longtable}{|p{0.6cm}|p{4.4cm}|p{2.8cm}|p{3.6cm}|p{2.6cm}|}
    \caption{Fallas reveladas por la evaluación. Las filas M son comportamientos del motor, no fallas del artefacto}
    \label{tab:res-fallas} \\
    \hline
    \textbf{Id} & \textbf{Falla} & \textbf{Detección} & \textbf{Efecto medido} & \textbf{Estado} \\
    \hline
    \endfirsthead
    \hline
    \textbf{Id} & \textbf{Falla} & \textbf{Detección} & \textbf{Efecto medido} & \textbf{Estado} \\
    \hline
    \endhead
    A1 & La personalización por currículo reemplazaba la consulta & Revisión frente a la documentación & Mismas ofertas para cualquier pregunta & Corregida (It.~1) \\
    \hline
    A2 & Las plantillas de \enquote{remoto} buscaban la palabra en la ubicación & Revisión de la agregación & 44 ofertas reconocidas frente a 682 & Corregida (It.~1) \\
    \hline
    A3 & El umbral de similitud de 0,5 no filtraba & Inspección de similitudes & Mínima 0,66 y mediana 0,75 en 20 consultas & Reemplazado por compuerta (It.~2) \\
    \hline
    A4 & Las restricciones de la consulta no llegaban a la búsqueda & Primera corrida & C@5 de 0,41 sin filtros frente a 0,95 con filtros & Corregida (It.~2) \\
    \hline
    A5 & El normalizador no convierte monedas ni montos por hora & Respuestas de agregación & Salario promedio del corpus de 544.127 USD & Mitigada (It.~2) \\
    \hline
    A6 & Respuestas en inglés a consultas en español & Primera corrida (PE2) & 21 de 25 respuestas & Corregida (It.~2): 4 de 24 \\
    \hline
    A7 & La plantilla de ofertas recientes se atendía por la ruta semántica & Corrida de prueba & Exactitud de la agregación 0,667 & Corregida: 0,733 \\
    \hline
    A8 & La palabra \textit{hiring} activa el patrón de recuperación & Corrida de prueba & 2 de 15 consultas de agregación mal enrutadas & Abierta \\
    \hline
    B1 & La consulta ordenaba por $1-d$ y no usaba el índice & \texttt{EXPLAIN ANALYZE} & 25,7~ms frente a 2,8~ms & Corregida (It.~1) \\
    \hline
    B2 & Índice IVFFlat creado sobre la tabla vacía & \texttt{EXPLAIN ANALYZE} & Equivalía a un recorrido completo & Reemplazado por HNSW (It.~1) \\
    \hline
    C1 & La ruta de agregación fallaba sin plantilla & Revisión del código & Error de ejecución & Corregida (It.~1) \\
    \hline
    C2 & Faltan plantillas: contar con una condición y contar empresas distintas & Corrida de prueba & Respuesta incorrecta con enrutamiento correcto & Abierta \\
    \hline
    D1 & El juez sin condiciones duras aceptaba ofertas que no cumplían la restricción & Revisión de juicios & La métrica no distinguía si un sistema respeta las restricciones & Corregida: juez v2 \\
    \hline
    D2 & Ventana de 4.096 \textit{tokens} del juez de RAGAS & Revisión de salidas del juez & Fidelidad truncada en 21 de 85 respuestas & Corregida: 16.384 \\
    \hline
    D3 & El juez local asigna fidelidad 0 a cifras presentes en el contexto & Revisión de casos extremos & PE2 no confiable con el juez local & Pendiente: juez final \\
    \hline
    D4 & Exportación limitada a 20 registros e ingesta no desactivable & Preparación del corpus & Habrían alterado el corpus congelado & Corregidas (It.~1) \\
    \hline
    D5 & El cálculo del kappa usaba la versión anterior del prompt del juez & Revisión del arnés & Habría comparado 269 de los 401 pares con veredictos del prompt anterior & Corregida \\
    \hline
    D6 & Los reintentos automáticos del cliente del generador contaban como latencia & Revisión de la latencia (PE3) & 59 de 90 generaciones con TTFT mayor que 2,5~s; mediana de 7,1~s & Corregida: reintentos desactivados \\
    \hline
    M1 & Un filtro selectivo en pgvector sin barrido iterativo devuelve listas incompletas & Benchmark (PE4) & Corpus real, filtro del 50\,\%: listas completas en el 22\,\% de las consultas & Evitado en el diseño (It.~1) \\
    \hline
    M2 & Con HNSW y pocas filas, el planificador prefiere el barrido secuencial & \texttt{EXPLAIN ANALYZE} & 24,8~ms frente a 2,0~ms con el índice forzado & Documentada; modo que fuerza el índice \\
    \hline
\end{longtable}
}

\subsection{Fallas silenciosas que cambian los resultados}

\begin{description}
    \item[A1. Personalización por currículo.] La documentación indicaba que el perfil del usuario se antepone a su consulta; la implementación construía la búsqueda solo con el perfil. Un usuario con currículo recibía las mismas ofertas fuera cual fuera su pregunta, y el modelo redactaba para cada pregunta una respuesta coherente con esas ofertas. Se corrigió para anteponer el perfil. \emph{Lección:} una ruta de personalización debe probarse variando la consulta y verificando que el resultado cambie.
    \item[A2. Ofertas remotas en la agregación.] Las plantillas que cuentan y listan ofertas remotas buscaban la palabra \textit{remote} en el texto de la ubicación y no usaban el atributo de modalidad que la ingesta obtiene de la fuente o del análisis estructurado. En el corpus de evaluación reconocían 44 ofertas remotas, frente a 682 marcadas con el atributo. El modelo informaba la cifra sin advertir nada. Se corrigió para usar el atributo, con el texto de la ubicación como respaldo. \emph{Lección:} cada plantilla de agregación debe contrastarse con una consulta de referencia escrita a mano, porque el modelo presenta cualquier cifra con la misma seguridad.
    \item[A3. Umbral de similitud.] El pipeline descartaba los fragmentos con similitud menor que 0,5. Con e5-base-v2, la similitud mínima entre la consulta y cualquier fragmento fue 0,66 en 20 consultas, y la mediana 0,75, de modo que el umbral nunca descartaba nada y toda consulta, incluso fuera de dominio, recibía cinco ofertas. La Iteración~2 lo reemplazó por una compuerta calibrada sobre el puntaje del reordenamiento (Sección~\ref{sec:art-iter-recuperacion}). \emph{Lección:} los umbrales de similitud dependen de la distribución de cada modelo y deben calibrarse con consultas dentro y fuera de dominio.
    \item[A4. Restricciones ignoradas.] Las consultas con señales de recuperación iban directamente a la búsqueda semántica y sus condiciones de salario, modalidad o tipo de contrato se ignoraban (Sección~\ref{sec:art-iter-recuperacion}). Las respuestas mencionaban ofertas que no cumplían lo pedido, sin advertirlo. Sin filtros, solo el 41\,\% de las cinco primeras ofertas de las consultas con restricciones las cumplía; con los filtros de la Iteración~2, el 95\,\% (Tabla~\ref{tab:res-restricciones}). \emph{Lección:} la calidad de la recuperación debe medirse también con criterios deterministas, como el cumplimiento de las restricciones, que no dependen del juez.
    \item[A5. Normalización de salarios.] El normalizador no convierte monedas ni algunos montos por hora o mensuales. Entre otros casos, una oferta de despacho registraba 86.528.000 dólares anuales, y el salario promedio del corpus que respondía la ruta de agregación era de 544.127 dólares. La Iteración~2 registró como no informados los salarios fuera del intervalo de 10.000 a 1.000.000 dólares anuales o con un mínimo mayor que el máximo; la regla afectó a 12 ofertas y el promedio pasó a 158.786 dólares. Es una mitigación: el salario de esas ofertas se pierde en lugar de convertirse. \emph{Lección:} las respuestas de agregación heredan los errores de normalización de los datos, y deben revisarse los valores extremos antes de evaluar.
    \item[A6. Idioma de la respuesta.] En la primera corrida, 21 de las 25 consultas en español se respondieron en inglés: sin una instrucción explícita, el modelo respondía en el idioma del contexto. La Iteración~2 agregó a los tres prompts la instrucción de responder en el idioma de la pregunta; 4 de 24 respuestas siguieron en inglés. \emph{Lección:} el idioma de la respuesta es una propiedad que debe especificarse y medirse, no suponerse.
    \item[A7. Ofertas recientes.] Ante \textit{\enquote{what are the most recent job postings}}, el clasificador devolvía la etiqueta de recuperación junto con la plantilla de ofertas recientes, y la consulta se atendía por la ruta semántica con una respuesta plausible pero no ordenada por fecha. Se corrigió con una regla de enrutamiento: las plantillas de estadísticas, valores distintos y ofertas recientes se atienden siempre por la agregación (Anexo~\ref{anx:clasificacion}). La Sección~\ref{sec:res-agregacion} reporta la exactitud antes y después. \emph{Lección:} la salida del clasificador debe validarse con reglas deterministas cuando la etiqueta y la plantilla se contradicen.
    \item[A8. La palabra \textit{hiring}.] El patrón de recuperación incluye la palabra \textit{hiring}, de modo que \textit{\enquote{list all companies hiring engineers}} y \textit{\enquote{how many different companies are hiring}} se enrutan como consultas mixtas. La falla se detectó en el conjunto de prueba y no se corrigió: corregirla habría ajustado el artefacto a esas consultas. Queda como trabajo futuro. \emph{Lección:} los patrones léxicos de enrutamiento deben probarse con consultas que usan las mismas palabras con otra intención.
\end{description}

\subsection{Fallas silenciosas que cambian el costo}

\begin{description}
    \item[B1. Consulta que no usaba el índice.] La búsqueda ordenaba por la expresión $1 - d(\mathbf{q}, \mathbf{v})$ en lugar de por la distancia, y el planificador de PostgreSQL no puede servir ese orden desde el índice. Los resultados eran correctos, porque el barrido completo es exacto, pero la búsqueda tardaba 25,7~ms en lugar de 2,8~ms (Tabla~\ref{tab:art-explain}). Se corrigió en la Iteración~1. \emph{Lección:} el plan de ejecución de cada consulta vectorial debe inspeccionarse; un índice que existe no garantiza que se use.
    \item[B2. Índice creado sobre la tabla vacía.] El índice IVFFlat se creó en la migración inicial, antes de cargar datos, de modo que todos los vectores quedaron en una sola lista y el índice equivalía a un recorrido completo. Se reemplazó por HNSW, que no requiere entrenamiento. \emph{Lección:} los índices que se entrenan sobre los datos deben construirse después de cargarlos.
\end{description}

\subsection{Fallas visibles}

\begin{description}
    \item[C1. Agregación sin plantilla.] Cuando los patrones clasificaban una consulta como agregación sin seleccionar ninguna plantilla, la ejecución fallaba con un error. Las pruebas de integración no lo detectaban porque simulaban el repositorio SQL. Se corrigió delegando la selección de la plantilla al clasificador basado en LLM, con la recuperación como alternativa, y las pruebas verifican ahora la plantilla seleccionada. \emph{Lección:} las pruebas que simulan el acceso a datos pueden ocultar fallas de integración.
    \item[C2. Plantillas faltantes.] El catálogo no tiene una plantilla que cuente las ofertas que cumplen una condición de salario, ni una que cuente empresas distintas. Ante \textit{\enquote{how many jobs pay more than 200k}}, el sistema elige correctamente la plantilla más cercana, que lista ofertas, y la respuesta no contiene el conteo pedido. Se detectó en el conjunto de prueba y queda abierta, por la misma razón que A8. \emph{Lección:} un catálogo cerrado de plantillas debe evaluarse por cobertura de las preguntas, no solo por exactitud del enrutamiento.
\end{description}

\subsection{Fallas del arnés de evaluación}

\begin{description}
    \item[D1. Juez sin condiciones duras.] La primera versión del prompt del juez calificaba como relevante una oferta del cargo correcto aunque no cumpliera el salario, la modalidad o el tipo de contrato pedidos. Con ese juez, la métrica no distinguía a un sistema que respeta las restricciones de uno que las ignora. La segunda versión trata esas condiciones como duras (Sección~\ref{sec:met-juicios}). Este cambio, junto con el del agrupamiento, explica que el nDCG@10 de la recuperación densa con e5-base-v2 pasara de 0,73 en la primera corrida a 0,68 en la definitiva sin que cambiara el sistema (Sección~\ref{sec:res-pe1}). \emph{Lección:} el prompt del juez debe expresar las condiciones que la métrica pretende medir.
    \item[D2. Ventana de contexto del juez de RAGAS.] Con la ventana por defecto de 4.096 \textit{tokens}, el contexto recuperado y las instrucciones no cabían, y la verificación de la fidelidad quedó truncada en 21 de las 85 respuestas. Todas se puntuaron de nuevo con una ventana de 16.384. \emph{Lección:} el arnés debe verificar que el juez recibe completo lo que juzga.
    \item[D3. Errores del juez local.] El juez local asignó fidelidad 0 a respuestas cuyas cifras y títulos aparecen literalmente en el contexto. Por eso la fidelidad medida con ese juez no se considera fiable, y se repetirá con el juez final y una verificación humana. \emph{Lección:} los puntajes extremos de un juez automático deben revisarse a mano antes de interpretarse.
    \item[D4. Exportación e ingesta.] El listado usado para exportar el catálogo estaba limitado a 20 registros, y la ingesta programada no podía desactivarse. Lo primero habría truncado el corpus exportado y lo segundo lo habría modificado durante la evaluación. Ambos se corrigieron antes de congelar el corpus. \emph{Lección:} la reproducibilidad exige poder congelar los datos y verificar su integridad.
    \item[D5. Versión del prompt en el cálculo del kappa.] El script que calcula el acuerdo entre el juez y los juicios humanos seleccionaba los veredictos de la primera versión del prompt, fijada en su código, mientras que la muestra se extrajo de los juicios de la segunda. Como la caché conserva los veredictos de ambas versiones, el cálculo no habría fallado: habría emparejado solo 269 de los 401 pares, con los veredictos del prompt anterior. Se detectó al preparar la elección del juez final, antes de calcular ningún kappa, y se corrigió haciendo explícita la versión del prompt. \emph{Lección:} cada artefacto del arnés que filtra resultados debe verificar que la cantidad de elementos emparejados es la esperada.
    \item[D6. Reintentos del cliente en la medición de la latencia.] La biblioteca cliente del proveedor del generador reintenta por su cuenta, antes que el arnés, las solicitudes rechazadas por límite de uso, y espera el tiempo que el proveedor indica. En una corrida de la latencia de la generación, esa espera se sumó al tiempo hasta el primer \textit{token}: 59 de las 90 generaciones superaron 2,5~s, con una mediana de 7,1~s, mientras que las generaciones sin espera tardaban 0,48~s. La corrida se descartó, y la latencia se mide ahora con los reintentos automáticos del cliente desactivados y una pausa de 35~s entre solicitudes (Sección~\ref{sec:met-metricas}). \emph{Lección:} la latencia de un servicio externo debe medirse con los reintentos del cliente desactivados, porque la espera que introducen no se distingue de la latencia del servicio.
\end{description}

\subsection{Comportamientos del motor}
\label{sec:res-motor}

La evaluación reveló además dos comportamientos de pgvector que no son defectos del artefacto, pero que tienen la misma naturaleza silenciosa: el motor devuelve resultados sin ningún aviso, aunque estén incompletos o se obtengan por un camino más lento. El diseño debe tenerlos en cuenta.

\begin{description}
    \item[M1. Filtros sin barrido iterativo en pgvector.] Con el índice HNSW, pgvector aplica el filtro sobre los candidatos que recorre el índice; si el filtro es selectivo, puede devolver menos de $k$ resultados sin ningún aviso. En el corpus real, con un filtro del 50\,\% y $ef_{search}=16$, solo el 22\,\% de las consultas obtuvo la lista completa y el \textit{recall} fue de 0,74; con barridos iterativos, todas la obtuvieron y el \textit{recall} subió a 0,93. El artefacto nunca lo sufrió, porque antes de la Iteración~2 los filtros no llegaban a la búsqueda vectorial, y el diseño de la Iteración~1 ya activaba esos barridos en toda búsqueda filtrada. \emph{Lección:} una búsqueda filtrada debe verificar que devuelve el número de resultados pedido.
    \item[M2. Preferencia del planificador.] Aun con la consulta corregida y un índice HNSW, el planificador de PostgreSQL elige el barrido secuencial a la escala del corpus real: 24,8~ms frente a 2,0~ms con el índice forzado. No es un error del artefacto sino un comportamiento del motor; la búsqueda admite un modo que fuerza el índice y la evaluación mide ambos. \emph{Lección:} la latencia de un motor vectorial depende de las decisiones de su planificador, que deben medirse a la escala real de uso.
\end{description}

\subsection{Síntesis}

De las 12 fallas del artefacto, 10 son silenciosas: ocho cambian los resultados y dos el costo. Ninguna producía un error visible ni una respuesta evidentemente incorrecta, porque el modelo de lenguaje convertía en texto fluido lo que recibía. Se detectaron porque el protocolo mide salidas intermedias y no solo la respuesta final: el plan de ejecución de la consulta, el número de resultados de una búsqueda filtrada, el cumplimiento de las restricciones, la plantilla elegida, el idioma de la respuesta y los veredictos del juez. Las seis fallas del arnés muestran que el instrumento de medición también falla en silencio y debe verificarse con el mismo cuidado que el artefacto. La Sección~\ref{sec:res-discusion} retoma estas observaciones.

\section{Discusión}
\label{sec:res-discusion}

\pendiente{Se redacta con la corrida final. Estructura prevista, por pregunta de evaluación:}
\begin{itemize}
    \item \pendiente{PE1: la recuperación densa supera a BM25, con la mayor ventaja en las consultas de desajuste de vocabulario, frente a la robustez de BM25 que reporta \citet{Thakur2021}; los filtros y el reordenamiento aportan mejoras significativas; la búsqueda híbrida no.}
    \item \pendiente{El resultado nulo de la búsqueda híbrida frente a la literatura: BM25 es más débil que la densa en este corpus, sobre todo en español, y RRF con $k = 60$ sin ajustar le da el mismo peso; \citet{Bruch2023} anticipan la sensibilidad de RRF y la ventaja de una combinación convexa.}
    \item \pendiente{PE2: fidelidad de las respuestas con el juez final y verificación humana.}
    \item \pendiente{PE3: el costo del reordenamiento frente a su aporte, y alternativas: menos candidatas, un modelo más pequeño o ejecutarlo solo cuando la primera etapa es incierta.}
    \item \pendiente{PE4: el motor no cambia la calidad a la escala real; con filtros selectivos y a gran escala Qdrant aventaja a pgvector, cuyo planificador elige el índice HNSW cuando la búsqueda exacta sería igual de rápida (\citealp{Sanca2024}).}
    \item \pendiente{Las fallas silenciosas: la fluidez del LLM enmascara defectos de recuperación, de datos y de enrutamiento; evaluar un sistema RAG exige verificar componente por componente y no solo la respuesta final (Sección~\ref{sec:res-fallas}; \citealp{Barnett2024}).}
    \item \pendiente{Limitaciones: juez, tamaño del corpus, un solo dominio e idioma de las ofertas, consultas redactadas por el autor.}
\end{itemize}
